\documentclass[12pt,a4paper]{report}
\usepackage{graphicx}
\usepackage{fancyhdr}
\usepackage{titlesec}
\usepackage{geometry}
\usepackage{placeins}
\usepackage{graphicx}
\usepackage{listings}
\usepackage{indentfirst}
\usepackage{xcolor}
\usepackage{float}
\usepackage{setspace}
\onehalfspacing

\lstdefinestyle{terminal}{
    backgroundcolor=\color{black},
    basicstyle=\ttfamily\color{white},
    frame=single,
    breaklines=true
}

\lstset{style=terminal}

\usepackage{fancyhdr}

% Configure fancyhdr
\pagestyle{fancy}
\fancyhf{} % Clear default header/footer
\fancyhead[L]{\leftmark} % Chapter name on the left
\fancyfoot[R]{\thepage} % Page number on the right
\renewcommand{\chaptermark}[1]{\markboth{\chaptername\ \thechapter: #1}{}}

% Remove header from TOC and Introduction
\fancypagestyle{plain}{ % Applies to chapter opening pages
    \fancyhf{} % Clear header/footer
    \fancyhead[L]{\leftmark} % Reintroduce chapter header
    \fancyfoot[R]{\thepage} % Page number on the right
}

% Apply an empty page style for TOC and Introduction
\fancypagestyle{noheader}{
    \fancyhf{} % Completely clear header/footer
    \fancyfoot[R]{\thepage} % Page number on the right
}

% Set margins
\geometry{
  a4paper,
  left=2.5cm,
  right=2.5cm,
  top=2.5cm,
  bottom=2.5cm,
  headsep=1.25cm,
  footskip=1.25cm
}

% Set paragraph indentation and spacing
\setlength{\parindent}{0.75cm}
\setlength{\parskip}{0cm}

% Define title style for chapters
\titleformat{\chapter}[block]
  {\normalfont\huge\bfseries}
  {Chapter \thechapter:}
  {1em} {}

\usepackage{fancyhdr}

% Configure header globally
\pagestyle{fancy}
\fancyhf{} % Clear default header/footer
\fancyhead[L]{\leftmark} % Chapter name on the left
\fancyfoot[R]{\thepage} % Page number on the right

% Ensure chapter names appear correctly
\renewcommand{\chaptermark}[1]{\markboth{\chaptername\ \thechapter: #1}{}}

\begin{document}\newpage
\pagestyle{empty}  % Remove headers and page numbers
\renewcommand{\headrulewidth}{0pt} % Remove horizontal line in header
\renewcommand{\footrulewidth}{0pt} % Remove footer line
\listoffigures
\thispagestyle{empty}  % Ensure this page has no formatting

\newpage
\pagestyle{empty}  % Apply empty style to Bibliography page
\renewcommand{\headrulewidth}{0pt}
\renewcommand{\footrulewidth}{0pt}
\renewcommand{\bibname}{Webography}

% Bibliography
\begin{thebibliography}{99}

\bibitem{rhel_security} Red Hat. \textit{Disk Partitioning Best Practices for Secure RHEL Deployments}.  
Available at: \url{https://docs.redhat.com/en/documentation/red_hat_enterprise_linux/9/html-single/security_hardening/index#Disk_partitioning_securing-rhel-during-installation}  

\bibitem{docker_docs} Docker Inc. \textit{Official Docker Documentation}.  
Available at: \url{https://docs.docker.com/}  

\bibitem{docker_logs} Docker Curriculum. \textit{Understanding Logs and Containerized Applications}.  
Available at: \url{https://docker-curriculum.com}  

\bibitem{datadog_aide} Datadog. \textit{Automating File Integrity Monitoring with AIDE and Cron}.  
Available at: \url{https://docs.datadoghq.com/security/default_rules/xccdf-org-ssgproject-content-rule-aide-periodic-cron-checking/}  

\bibitem{hands_on_docker} LabEx. \textit{Hands-on Docker Security and Monitoring Course}.  
Available at: \url{https://labex.io}  

\bibitem{tryhackme_containers} TryHackMe. \textit{Practical Container Security Challenges and Linux Hardening}.  
Available at: \url{https://tryhackme.com}  

\end{thebibliography}
\thispagestyle{empty} % Ensure Bibliography has no formatting

\newpage
% Acknowledgments Section
\pagestyle{empty}  % Remove header/footer for this page
\renewcommand{\headrulewidth}{0pt}
\renewcommand{\footrulewidth}{0pt}
\chapter*{Acknowledgments}

I would like to express my sincere gratitude to \textbf{Ooredoo} for providing me with the opportunity to undertake this internship. The experience has been invaluable in expanding my knowledge and practical expertise in container security and system hardening.

A special thanks to \textbf{Ms. Imen Aloui}, my supervisor, for her unwavering guidance, insightful feedback, and continuous support throughout the duration of this internship. Her expertise and encouragement have played a crucial role in shaping my understanding of the subject matter and overcoming various challenges.

This experience has not only enhanced my technical skills but has also provided me with a deeper appreciation for security best practices in modern IT infrastructures. I am truly grateful for the opportunity and the knowledge I have gained through this journey.

\thispagestyle{empty} % Ensure Acknowledgments page has no formatting

\newpage
\pagestyle{fancy} % Restore normal headers/footers after these special sections


\pagestyle{empty} % Remove headers/footers from TOC
\renewcommand{\headrulewidth}{0pt} % Remove horizontal line
\tableofcontents
\thispagestyle{empty} % Ensure first TOC page has no numbering

% Ensure subsequent TOC pages also have no numbering
\addtocontents{toc}{\protect\thispagestyle{empty}}


\newpage

% Define a custom page style for Introduction (No header, keep page number)
\fancypagestyle{noheaderplain}{
    \fancyhf{} % Clear header and footer
    \fancyfoot[R]{\thepage} % Page number on the right
    \renewcommand{\headrulewidth}{0pt} % Remove top horizontal line
}

\newpage
\pagestyle{fancy} % Restore headers after TOC
\renewcommand{\headrulewidth}{0.4pt} % Restore header line for the rest of the document
\fancyfoot[R]{\thepage} % Page number on the right
\pagenumbering{arabic}
\setcounter{page}{1}

% Introduction (No header, but keep numbering)
\chapter*{Introduction}
\addcontentsline{toc}{chapter}{Introduction} % Add to TOC without numbering
\thispagestyle{noheaderplain} % Use the custom style for Introduction




The growing adoption of containerized environments has brought new security issues, and strong security practices must be applied to contemporary IT infrastructures. This report analyzes some of the most important security principles applicable to containerization, including compliance, system hardening, and architectural differences between Docker and Podman. It involves applying industry best practices to secure containerized workloads, such as conducting vulnerability scanning and compliance scanning with the use of tools such as Trivy, OpenSCAP, and AIDE. Docker and Podman are also compared against each other with emphasis on the security benefits of Podman's rootless mode of execution. Monitoring and logging systems such as Prometheus, Grafana, and Fluent Bit are also incorporated to add system observability and threat detection.

Through the integration of theoretical review and practical application, this report offers an understanding of enhancing container security, reducing vulnerabilities, and resolving legacy system issues. The project emphasizes the need for proactive security measures to ensure the resilience of containerized infrastructure.


\newpage

% Chapter 1 - Company Presentation
\setcounter{chapter}{0} % Ensure Chapter 1 starts at 1
\chapter{Presentation of the Host Organization and Security Division}

\section{Host Organization Overview}

\subsection{History}
Ooredoo Tunisia, originally established as Tunisiana on May 11, 2002, was the first private mobile operator in Tunisia. It launched commercial activities on December 27, 2002, and within six months, its network coverage had reached 60\% of the population. By 2005, this expanded to 99\%.

Over the years, the company underwent several ownership transitions. In 2010, Wataniya Telecom, a subsidiary of Qatar Telecom (now Ooredoo), acquired a majority stake. Following the Tunisian revolution, the government took temporary control of a portion of these shares before selling a 15\% stake to Qatar Telecom in 2012. The company officially rebranded as Ooredoo Tunisia on April 24, 2014.

\subsection{Company Overview}
As part of the Ooredoo Group, which operates in 10 countries and serves over 130 million customers, Ooredoo Tunisia has become the country’s largest telecommunications provider. It offers mobile, fixed, broadband internet, and enterprise services.

Structured as a private limited company, Ooredoo Tunisia has a capital of 359 million Tunisian dinars and serves over 7.5 million subscribers. The company collaborates with more than 5,000 suppliers, 95\% of whom are Tunisian enterprises. Its headquarters is located in the "Z\'enith" building, Les Jardins du Lac, Tunis.


\begin{figure}[H]
    \centering
    \begin{minipage}{0.35\textwidth}
        \centering
        \includegraphics[width=0.8\linewidth]{images/tunisiana.jpg}
        \caption{Tunisiana Logo (2006-2014)}
    \end{minipage}
    \hfill
    \begin{minipage}{0.35\textwidth}
        \centering
        \includegraphics[width=0.8\linewidth]{images/Ooredoo.png}
        \caption{Ooredoo Logo (2014-Present)}
        \label{fig:logos} % Moved inside to prevent referencing issues
    \end{minipage}
\end{figure}



\subsection{Corporate Structure}
Ooredoo Tunisia is composed of several key departments that ensure smooth operations:

\begin{itemize}
    \item \textbf{Administrative and Financial Department (DAF)}: Manages financial resources and internal controls.
    \item \textbf{Legal and Regulatory Affairs Department (DJR)}: Ensures compliance with legal and regulatory requirements.
    \item \textbf{Customer Service Department (DSC)}: Handles customer support and retention strategies.
    \item \textbf{Technical Department (DT)}: Oversees network design, deployment, maintenance, and optimization.
    \item \textbf{Commercial Department (DC)}: Develops pricing strategies and manages distribution channels.
    \item \textbf{Information Systems Department (DSI)}: Implements and maintains IT solutions supporting operations.
    \item \textbf{Quality and Programs Department (DQP)}: Conducts quality assessments and audits.
    \item \textbf{Human Resources Department (DRH)}: Manages recruitment, training, and employee relations.
\end{itemize}

Among these, the \textbf{Technical Department} plays a crucial role in maintaining Ooredoo’s telecommunications infrastructure, ensuring network security, and implementing cybersecurity measures.

\section{Internship Department: Security Division}
As part of the internship, work was conducted within the Security Division of the Technical Department. This division is responsible for safeguarding Ooredoo Tunisia's infrastructure against cyber threats and ensuring compliance with security best practices.

The primary responsibilities of the Security Division include:
\begin{itemize}
\item Implementing and managing security protocols to protect company assets.
\item Conducting risk assessments and vulnerability analyses.
\item Ensuring compliance with national and international cybersecurity regulations.
\item Developing incident response strategies and disaster recovery plans.
\item Enhancing security awareness across the organization.
\end{itemize}

The internship primarily focused on a self-directed learning project related to security. This involved exploring various cybersecurity concepts, best practices, and tools applicable to securing containerized environments. The knowledge gained through this process will be further detailed in subsequent chapters.

\section{Conclusion}
This chapter provided an overview of Ooredoo Tunisia, detailing its history, corporate structure, and the Security Division’s role within the company. Understanding the organizational framework was crucial in contextualizing my internship experience. The next chapter will focus on the technical aspects and specific security measures implemented during my time at Ooredoo Tunisia.



% Ensure page numbering starts at 1 for the first chapter page


\newpage
\chapter{Core Concepts and Principles}


\section{Introduction}
Containerization technologies have redefined application deployment paradigms by enabling lightweight, isolated environments that optimize resource utilization and portability. While Docker remains the industry standard for container orchestration, emerging alternatives like Podman introduce architectural innovations addressing critical security limitations. This chapter systematically analyzes core containerization principles, beginning with Docker’s runtime architecture and its divergence from traditional virtual machines. Subsequent sections evaluate security challenges inherent to containerized workflows, including vulnerability management, compliance frameworks, and hardening methodologies. By juxtaposing Docker with Podman’s daemonless design, this chapter establishes foundational knowledge for implementing secure containerization strategies.

\section{ Docker}
Docker is an open-source project that automates the deployment of software applications inside containers. It provides an additional layer of abstraction and automation of OS-level virtualization in Linux, making it easier to develop, ship, and run applications.

\subsection{Podman: An Alternative to Docker}
Podman offers container management without the need for a daemon. It provides improved security and maintains compatibility with Docker images, making it a viable alternative for container orchestration.

\textbf{What is a Daemon?}
A daemon is a background process that runs continuously, managing system tasks without direct user interaction. In the context of containerization, Docker relies on a central daemon (dockerd) to handle container lifecycle operations. In contrast, Podman operates in a daemonless manner, allowing users to run containers as independent processes, enhancing security and reducing the risk of a single point of failure.
\newpage
\section{Definition  of Containerization}
Containers offer a logical packaging mechanism in which applications can be abstracted from the environment in which they actually run. This decoupling allows container-based applications to be deployed easily and consistently, regardless of whether the target environment is a private data center, the public cloud, or even a developer’s personal laptop. This gives developers the ability to create predictable environments that are isolated from the rest of the applications and can be run anywhere.

From an operations standpoint, apart from portability, containers also give more granular control over resources, giving your infrastructure improved efficiency which can result in better utilization of your computer resources.
\section{Difference between Containers and Virtual Machines}
Containers and virtual machines (VMs) both provide isolated environments for running applications, but they differ in their architecture. Containers share the host OS kernel and isolate processes, whereas virtual machines run separate OS instances on a hypervisor.
\begin{figure}[h]
    \centering
    \includegraphics[width=0.8\textwidth]{images/vms-vs-containers.png}
    \caption{Difference between Containers and Virtual Machines}
\end{figure}

% Requires: \usepackage{graphicx}
\begin{table}[h]
    \centering
    \resizebox{\textwidth}{!}{%
    \begin{tabular}{|l|p{6cm}|p{6cm}|} \hline 
    \hliney
    \textbf{Differences} & \textbf{Docker} & \textbf{Virtual Machine} \\ \hline
    \textbf{Operating System} & 
    Docker is a container-based model where containers are software packages used for executing an application on any operating system. In Docker, the containers share the host OS kernel. Here, multiple workloads can run on a single OS. & 
    It is not a container-based model; they use user space along with the kernel space of an OS. It does not share the host kernel. Each workload needs a complete OS or hypervisor. \\ \hline
    \textbf{Performance} & 
    Docker containers result in high performance as they use the same operating system with no additional software (like hypervisor). Docker containers can start up quickly and result in less boot-up time. & 
    Since VM uses a separate OS; it causes more resources to be used. Virtual machines don’t start quickly and lead to poor performance. \\ \hline
    \textbf{Portability} & 
    With docker containers, users can create an application and store it into a container image. Then, he/she can run it across any host environment. Docker container is smaller than VMs, because of which the process of transferring files on the host’s filesystem is easier. & 
    It has known portability issues. VMs don’t have a central hub and it requires more memory space to store data. While transferring files, VMs should have a copy of the OS and its dependencies because of which image size is increased and becomes a tedious process to share data. \\ \hline
    \textbf{Speed} & 
    The application in Docker containers starts with no delay since the OS is already up and running. These containers were basically designed to save time in the deployment process of an application. & 
    It takes a much longer time than it takes for a container to run applications. To deploy a single application, Virtual Machines need to start the entire OS, which would cause a full boot process. \\ \hline
    \end{tabular}%
    }
    \caption{Comparison between Docker and Virtual Machine}
    \label{tab:docker_vs_vm}
\end{table}
\FloatBarrier


\section{Docker Architecture}
\subsection{Components}
\begin{itemize}
    \item \textbf{Docker Engine:} The fundamental component responsible for managing and running containers.  
    \item \textbf{Docker Images:} Pre-configured templates used to instantiate containers.  
    \item \textbf{Docker Containers:} Lightweight, executable instances derived from Docker images.  
    \item \textbf{Docker Registry:} A centralized repository to store and distribute Docker images.  
\item \textbf{Docker daemon:} A background service that listens to Docker API requests, manages container lifecycles, builds images, and handles networking and storage.  

\end{itemize}

\begin{figure}[h]
    \centering
    \includegraphics[width=1\textwidth]{images/docker-arch.png}
    \caption{Docker Architecture}
\end{figure}

\subsection{How Containers Work Internally}
Docker operates through several key components that work together to manage containerized environments efficiently. The Docker Engine is responsible for handling container operations and interacts with the Docker Daemon a background service that listens for API requests, builds images, manages networking, and controls container lifecycles. Docker Images serve as immutable templates containing application code and dependencies, utilizing a layered filesystem (UnionFS to optimize storage and reduce redundancy. Docker Containers are runtime instances of images, isolated from the host system using namespaces for process and network separation, and cgroups to allocate system resources. The Docker Registry functions as a centralized storage system, enabling the distribution, versioning, and retrieval of images for deployment.

\section{Security Risks in Containerized Environments}

Containerized environments present various security challenges that must be addressed to ensure a secure deployment. The following are key security risks associated with containerized applications:

\begin{itemize}
    \item \textbf{Insecure Container Images:}  
    Container images may contain known vulnerabilities, posing a risk to the entire infrastructure. The use of outdated or unverified images increases the likelihood of exploitation. Regular vulnerability scanning and the use of trusted image sources are essential to mitigate this risk.

    \item \textbf{Misconfigured Container Runtimes:}  
    Improper configuration of container runtimes can weaken isolation mechanisms, increasing the risk of unauthorized access and privilege escalation. Ensuring secure runtime configurations and enforcing strict isolation policies are crucial to maintaining container security.

    \item \textbf{Weak or Outdated Container Dependencies:}  
    Containers often rely on external dependencies, including libraries and frameworks, which may contain security vulnerabilities. The presence of outdated or unpatched dependencies increases the risk of exploitation. A structured update and patch management strategy is necessary to minimize such risks.

    \item \textbf{Insecure API Interfaces:}  
    Unsecured container APIs expose systems to unauthorized access and potential exploitation. It is essential to enforce strong authentication, access controls, and monitoring mechanisms to protect exposed API endpoints.

    \item \textbf{Insider Threats:}  
    Weak access control mechanisms can lead to insider threats, where unauthorized individuals exploit excessive privileges to access sensitive data. Implementing role-based access controls (RBAC) and continuous monitoring helps mitigate these risks.

    \item \textbf{Data Breaches and Leakage:}  
    Improper handling of sensitive data within containerized environments can result in data breaches. Attackers may exploit misconfigurations to gain access to confidential information. Employing strong encryption, secure secret management practices, and strict access controls are essential for protecting sensitive data.

    \item \textbf{Container Breakouts:}  
    Exploiting vulnerabilities in the host operating system or kernel may allow an attacker to escape container isolation and gain control over the underlying system. Regular patching of the host OS, kernel, and container runtime is crucial to preventing such attacks.

    \item \textbf{Insecure Container Registries:}  
    The use of untrusted or poorly secured container registries increases the risk of deploying malicious or tampered images. Organizations should rely on reputable registries, enforce strict image validation policies, and implement signed container images to ensure integrity and authenticity.
\end{itemize}
\section{Compliance vs. Hardening in Container Security}

In container security, compliance and hardening are two fundamental practices that contribute to a secure environment. While closely related, they address different aspects of security. Compliance ensures that containerized systems adhere to industry standards, regulatory requirements, and security policies, while hardening focuses on minimizing vulnerabilities by applying security best practices, reducing attack surfaces, and strengthening system configurations. A clear understanding of both approaches is essential for establishing a resilient security framework.  

\subsection{Compliance}

Compliance refers to ensuring that containerized environments meet specific regulatory requirements, industry standards, or security benchmarks. The goal of compliance is to verify that the infrastructure adheres to established security guidelines and legal requirements. Compliance frameworks, like CIS Benchmark, outline a set of security best practices that organizations must follow. 


For containerized environments, compliance ensures:
\begin{itemize}
    \item Adherence to security best practices outlined by authoritative bodies.
    \item Validation of security configurations to meet specific regulatory standards.
    \item Ongoing audits and assessments to ensure continued adherence to compliance policies.
\end{itemize}

Tools such as Docker Bench for Security and OpenSCAP are commonly used to assess the security of containerized environments against recognized compliance standards.

\subsection{Hardening}

Hardening, on the other hand, involves tightening security measures to reduce the attack surface and minimize vulnerabilities within the containerized environment. The goal of hardening is to proactively configure and secure systems, containers, and their dependencies to withstand attacks, even in the absence of specific regulatory requirements.

Hardening includes:
\begin{itemize}
    \item Disabling unnecessary services and ports within containers to reduce the attack surface.
    \item Implementing strong authentication and access control mechanisms to limit privileged access.
    \item Regularly updating containers and their dependencies to patch known vulnerabilities.
    \item Configuring container runtimes securely to prevent unauthorized access and exploitations.
\end{itemize}

Hardening focuses on reducing the risk of exploitation, regardless of whether or not specific compliance standards are met. Tools like Kube-bench, Trivy, and Falco assist with detecting configuration flaws, vulnerabilities, and suspicious behavior within the container runtime and ecosystem.

\subsection{Key Differences}

While both compliance and hardening are essential for securing containerized environments, the primary difference lies in their objectives:
\begin{itemize}
    \item Compliance focuses on ensuring adherence to external standards, regulations, and guidelines.
    \item Hardening emphasizes strengthening the container environment by proactively addressing security risks, often beyond the scope of compliance requirements.
\end{itemize}
Both approaches are complementary, as achieving compliance often involves following a hardened security posture, and hardening practices benefit from being aligned with compliance standards.
\section{Differences Between Podman and Docker}

While Docker and Podman are both containerization tools designed to manage and orchestrate containers, they differ significantly in their architecture and security models. Understanding these differences is crucial for selecting the appropriate tool based on specific operational and security requirements.

\subsection{Architectural Differences}
The primary distinction between Docker and Podman lies in their architecture. Docker relies on a client-server model where the Docker CLI communicates with a background daemon (\texttt{dockerd}) responsible for managing containers. This daemon runs with root privileges, which introduces potential security risks if compromised. In contrast, Podman employs a daemonless architecture, where containers are managed directly by the Podman CLI without requiring a central daemon. This approach allows Podman to run containers as non-root users, thereby reducing the attack surface and enhancing security.

\subsection{Rootless Containers}
Podman's support for rootless containers is a significant security advantage. By allowing containers to execute without root privileges, Podman minimizes the risk of privilege escalation attacks. Docker, while offering rootless mode as an experimental feature, traditionally requires root access for daemon operations, which can pose security concerns in shared or untrusted environments.

\subsection{Compatibility and CLI}
Podman maintains compatibility with Docker's CLI, enabling users to transition seamlessly by using analogous commands (e.g., \texttt{podman} instead of \texttt{docker}). This compatibility extends to Docker images, as Podman can pull and run images from Docker registries like Docker Hub. However, Podman does not support Docker Swarm, instead relying on Kubernetes for orchestration.

\subsection{Security Implications}
The absence of a daemon in Podman eliminates a common attack vector present in Docker. Since Podman does not require a continuously running daemon, it reduces the risk of daemon-related vulnerabilities. Additionally, Podman's integration with systemd for managing containers as services enhances security through established process supervision and logging mechanisms.
\section{Conclusion}
This chapter has elucidated the architectural and operational principles underpinning modern containerization technologies. Key findings include:

Architectural Efficiency: Containers provide superior performance over virtual machines through kernel sharing, though this introduces shared-resource security risks.

Security Trade-offs: Docker’s daemon-based model simplifies orchestration but expands attack surfaces, while Podman’s rootless, daemonless approach enhances security at the cost of ecosystem maturity.

Risk Mitigation: Effective container security necessitates layered strategies combining image scanning, runtime hardening, and compliance auditing.

The analysis underscores that containerization’s benefits are inextricably linked to proactive security practices. As organizations adopt these technologies, understanding the Docker-Podman dichotomy and implementing rigorous hardening protocols become critical success factors. Subsequent chapters build upon this foundation to demonstrate practical security implementations.


\newpage
\newpage\chapter{Securing Containerized Environments}
\section{Introduction}

Securing containerized environments has become a critical aspect of modern infrastructure, particularly with the increasing adoption of containers in enterprise deployments. This chapter documents the process undertaken to explore container security, beginning with foundational containerization concepts and progressing towards advanced hardening techniques. Through hands-on experimentation with Docker and Podman, as well as the application of security tools like Trivy, OpenSCAP, and AIDE, a comprehensive understanding of securing containerized workloads was developed. The journey also involved overcoming challenges posed by legacy environments, reinforcing the importance of proactive security measures.
\section{Initial Exploration of Containerization}
At the outset of the internship, a foundational exploration of containerization was conducted to develop proficiency in Docker and containerized environments. This phase involved executing fundamental Docker commands, managing container lifecycles, and analyzing the architecture of container images.

To expand my understanding of containerization beyond Docker, I transitioned to using \textbf{Podman}, an alternative container engine that offers improved security features, such as rootless execution and better integration with SELinux. This switch allowed me to compare the behavior and security aspects of both technologies.

To practically apply these concepts, a basic web application was containerized and deployed. This deployment served as a test case for log analysis, enabling the detection of potential security threats. The environment was defined through a Dockerfile, automating the setup and ensuring consistency in container orchestration.
\newpage
\begin{lstlisting}[language=Dockerfile, caption={Dockerfile for the Log Analysis Web Application}, label={lst:dockerfile}, captionpos=b]
FROM ubuntu:22.04

RUN apt-get update && apt-get install -y python3 python3-pip && rm -rf /var/lib/apt/lists/*

WORKDIR /app

COPY requirements.txt . && pip3 install --no-cache-dir -r requirements.txt

COPY . .

EXPOSE 5000

CMD ["python3", "app.py"]
\end{lstlisting}

Following the containerization process, the application was deployed using Docker. The following commands were executed to build and run the container in detached mode, mapping port 8080 of the host system to port 5000 of the container:

\begin{lstlisting}
$ docker build -t web-loganalyzer .
$ docker run -d -p 8080:5000 web-loganalyzer
\end{lstlisting}

\begin{figure}[h]
    \centering
    \includegraphics[width=0.9\textwidth]{images/loganalyser.png}
    \caption{Deployed Web Application for Log Analysis}
    \label{fig:docker_commands}
\end{figure}

\section{Security Hardening of RHEL 9 Container Image}
Building upon initial containerization proficiency, the next phase focused on enhancing the security of a container image based on Red Hat Enterprise Linux 9 (RHEL 9). The hardening process encompassed: All these hardening measures were subsequently applied to a RHEL 7.4 container image as well, ensuring consistent security practices across different versions of the operating system. In the following sections, we will compare the results of applying the same hardening measures on RHEL 9 and RHEL 7 to assess the impact and effectiveness of the security enhancements on both versions.

\begin{itemize}
    \item Switching from Docker to \textbf{Podman} to evaluate its security benefits, including rootless execution .
    \item Performing an initial vulnerability scan using \textbf{Trivy} to identify security risks before applying hardening measures. The scan results revealed several medium and low vulnerabilities, particularly related to outdated packages and misconfigurations.
\end{itemize}

\begin{lstlisting}
$ trivy image rhel9:9.5
\end{lstlisting}


The scan report highlighted multiple vulnerabilities, necessitating targeted security improvements.
\begin{figure}[H]
    \centering
    \includegraphics[width=\textwidth]{images/trivyscan.png}
    \caption{First Trivy scan on rhel9}
    \label{fig:security_improvements}
\end{figure}

\begin{itemize}
    \item Removing privilege escalation when running the container to restrict unauthorized privilege elevation.
    \item Dropping unnecessary kernel capabilities to minimize potential attack vectors by using

    \item Implementing disk partitioning best practices by creating separate partitions for \texttt{/boot}, \texttt{/home}, \texttt{/tmp}, and \texttt{/var/tmp} to ensure data protection and prevent system instability.

\begin{lstlisting}
$ podman  run -t --security-opt=no-new-privileges -it \
  --name my-secure-container \
  --cap-drop all \
  -v $(pwd)/boot:/boot \
  -v $(pwd)/root:/mnt/root \
  -v $(pwd)/home:/home \
  -v $(pwd)/tmp:/tmp \
  -v $(pwd)/vartmp:/var/tmp \
  harden:9.5
\end{lstlisting}
    \item Using \textbf{OpenSCAP (OSCAP)} inside the container to validate security compliance and hardening effectiveness.
\end{itemize}
The OSCAP assessment provided a compliance report against predefined security benchmarks:
\begin{lstlisting}
$ sudo oscap xccdf eval --profile xccdf_org.ssgproject.content_profile_cis \
--remediate /usr/share/xml/scap/ssg/content/ssg-rhel9-xccdf.xml
\end{lstlisting}

The results indicated security gaps that were subsequently addressed through configuration adjustments. After remediation, subsequent scans showed significant improvement in the security score.
\begin{figure}[H]
    \centering
    \includegraphics[width=\textwidth]{images/OSCAPreport.png}
    \caption{OSCAP scan results after remediations}
    \label{fig:security_improvements}
\end{figure}



\begin{itemize}
    \item Setting up periodic system integrity checks using \textbf{AIDE} (Advanced Intrusion Detection Environment) combined with \textbf{Cronie} to automate file integrity verification.
\end{itemize}
\begin{lstlisting}
$ aide --init
$ mv /var/lib/aide/aide.db.new.gz /var/lib/aide/aide.db.gz
$ crontab -e
# Add the following line for periodic integrity checks
0 3 * * * /usr/sbin/aide --check
\end{lstlisting}

This phase provided essential insights into enterprise-grade container security measures and compliance requirements.
\section{Challenges and Key Insights}
Throughout the internship, several challenges were encountered, particularly with RHEL 7:
\begin{itemize}
    \item Many updates were impossible to resolve due to the lack of active repositories, as RHEL 7 is out of update support.
    \item Installing \textbf{dnf} was problematic, preventing the finalization of \textbf{OSCAP} and \textbf{AIDE} installation.
    \item Updates and package management were difficult, requiring manual intervention to find alternative package sources or workarounds.
    \item Issues with the subscription manager in the Dockerfile: The subscription manager did not provide access to repositories within the Dockerfile, necessitating manual updates and repository access within the container itself.
\end{itemize}

These obstacles highlighted the importance of maintaining up-to-date systems and the challenges of working with legacy software. Additionally, the manual intervention required for the subscription manager emphasized the need for streamlined processes in containerized environments to ensure efficient and automated updates.


\section{Implementation of Monitoring and Logging Mechanisms}
As part of the security hardening process, centralized monitoring and logging solutions were deployed to ensure real-time security oversight. The following tools were integrated:

\begin{itemize}
    \item \textbf{Prometheus and Grafana}: Used for real-time monitoring of container performance and security metrics.
    \item \textbf{Fluent Bit}: Implemented for efficient log collection and aggregation, enhancing observability of security events.
\end{itemize}

These tools were deployed using Docker and configured to operate on the same network for seamless communication:

\begin{lstlisting}
$ docker network create monitoring-net
$ docker run -d --name prometheus --network monitoring-net prom/prometheus
$ docker run -d --name grafana --network monitoring-net grafana/grafana
$ docker run -d --name fluentbit --network monitoring-net fluent/fluent-bit
\end{lstlisting}

Each container was assigned an IP within the network, ensuring proper connectivity and allowing web interfaces to be accessible for monitoring. The same setup was applied to a Podman-managed RHEL 9 container to evaluate its compatibility with the monitoring tools.
\newpage
\subsection{Conclusion}
The progression from basic containerization to  security hardening illustrated the complexities of securing containerized environments. The transition from Docker to Podman, vulnerability scanning, privilege restriction, and compliance validation were essential components of this journey. Additionally, the challenges encountered, particularly with RHEL 7, reinforced the need for maintaining updated software and efficient security management processes. By integrating monitoring and logging mechanisms, such as Prometheus, Grafana, and Fluent Bit, continuous oversight of container security was established. These combined efforts contribute to a more resilient and secure containerized infrastructure.
\chapter*{General conclusion}
\addcontentsline{toc}{chapter}{General conclusion} % Add to TOC without numbering
\markboth{General conclusion}{} % Show "Introduction" in the header
\thispagestyle{plain}
Ensuring the security of containerized environments requires a comprehensive approach that balances compliance, hardening, and continuous monitoring. This report has highlighted the critical aspects of container security, including the advantages of using Podman over Docker, the role of vulnerability assessment tools, and the necessity of implementing strong access controls.

A key challenge in securing containerized workloads lies in maintaining up-to-date container images and enforcing strict version control. While keeping images updated is essential to mitigate known vulnerabilities, the process is often complex due to dependency management, compatibility issues, and the risk of introducing instability into production environments. Additionally, working with legacy systems further complicates update procedures, as outdated software may lack active support or require extensive manual intervention to apply security patches.

Despite these challenges, adopting a structured security approach that integrates regular vulnerability scanning, compliance validation, and logging mechanisms significantly enhances the resilience of containerized infrastructures. The findings presented in this report underscore the importance of continuous security efforts, emphasizing that proactive measures are necessary to safeguard modern IT environments against evolving threats.


\end{document}
